\documentclass[11pt,a4paper]{article}

\usepackage[utf8]{inputenc}

\usepackage[T1]{fontenc}

\usepackage[polish]{babel}

\usepackage{lmodern}

\usepackage{geometry}

\geometry{margin=27mm}

\usepackage[hidelinks]{hyperref}

\title{Otwarty audyt błędów wnioskowania odwrotnego w neurokognicji}

\author{Mateusz Pomianek}

\date{Szkic naukowy, wersja 1.0.0 | 28 listopada 2026}

\begin{document}

\maketitle

\begin{abstract}

Proponowana miara główna to różnica alf Krippendorffa dla dwóch instrukcji kodowania.

\end{abstract}

\section{Wprowadzenie}

Wnioskowanie odwrotne zależy od swoistości wskaźnika i częstości bazowej procesu (Poldrack, 2006, 2011). Rubryka audytu powinna wymagać jawnego powiązania twierdzenia z tymi przesłankami.

\section{Cel i pytania badawcze}

Hipotezy robocze wymagają zatwierdzenia przed rozpoczęciem właściwej weryfikacji.

\paragraph{H1.} Zaślepieni koderzy korzystający z rubryki, która oddziela swoistość, częstość bazową i zakres wniosku, uzyskają wyższą alfę Krippendorffa dla klasyfikacji rodzaju błędu niż koderzy używający krótkiej instrukcji ogólnej.

\section{Materiał i metody}

Pierwszorzędowy wskaźnik: różnica alf Krippendorffa dla dwóch instrukcji kodowania.

Koderzy otrzymają zrównoważone, niepowtarzające się zestawy fragmentów z równoległymi formami i losowym przydziałem instrukcji. Przedziały dla różnicy zgodności będą uwzględniać grupowanie fragmentów według publikacji. Klasyfikacja przypadków kontrolnych pozwoli odróżnić zgodność od trafności kodowania.

\section{Plan analizy}

Zamrożenie estymandu, reguł oceny jakości i obsługi braków nastąpi przed dostępem do danych wynikowych. Raport obejmie oszacowanie kontrastu głównego z przedziałem niepewności i pełne mianowniki. Jeśli przewidziano hipotezy drugorzędowe, przed analizą należy ustalić sposób kontroli wielokrotności.

\section{Wyniki}

W niniejszym pliku nie umieszczono wyników empirycznych, statystyk ani nieprzeprowadzonych analiz. Sekcja wymaga uzupełnienia po uzyskaniu i weryfikacji materiału dowodowego.

\section{Dyskusja i ograniczenia}

Ograniczenia obejmują możliwe błędy pomiaru, dobór próby, alternatywne mechanizmy i ograniczoną możliwość uogólniania. Korelacja nie dowodzi przyczynowości; poprawność techniczna nie jest dowodem znaczenia klinicznego.

\section{Wnioski}

Hipotezy i sposób pomiaru stanowią podstawę oceny wykonalności oraz przygotowania właściwego protokołu.

\section*{Bibliografia}

{\setlength{\parindent}{0pt}
\setlength{\parskip}{5pt}
\hangindent=1.27cm
Poldrack, R. A. (2006). Can cognitive processes be inferred from neuroimaging data? \textit{Trends in Cognitive Sciences, 10}(2), 59--63. \url{https://doi.org/10.1016/j.tics.2005.12.004}\par
\hangindent=1.27cm
Poldrack, R. A. (2011). Inferring mental states from neuroimaging data: From reverse inference to large-scale decoding. \textit{Neuron, 72}(5), 692--697. \url{https://doi.org/10.1016/j.neuron.2011.11.001}\par
}

\end{document}
